\documentclass[14pt,a4paper]{extarticle}

% XeLaTeX รองรับ Unicode และอ่านฟอนต์ภาษาไทยจากไฟล์ในโครงการโดยตรง
\usepackage{fontspec}
\usepackage{polyglossia}
\setdefaultlanguage{thai}
\setotherlanguage{english}
% กำหนดกฎการตัดบรรทัดและระยะยืดระหว่างคำสำหรับข้อความภาษาไทย
\XeTeXlinebreaklocale "th"
\XeTeXlinebreakskip=0pt plus 1pt

\setmainfont[
  Path=fonts/,
  UprightFont={THSarabunNew.ttf},
  BoldFont={THSarabunNew Bold.ttf},
  ItalicFont={THSarabunNew Italic.ttf},
  BoldItalicFont={THSarabunNew BoldItalic.ttf}
]{THSarabunNew}

\usepackage{geometry}
\geometry{margin=2.54cm}

\title{แม่แบบเอกสารภาษาไทยด้วย \LaTeX}
\author{CRU LaTeX Template}
\date{\today}

\begin{document}

\maketitle

\section{ทดสอบภาษาไทย}
เอกสารนี้คอมไพล์ด้วย XeLaTeX ภายใน Docker และใช้ฟอนต์ TH Sarabun New
ที่เก็บอยู่ในโครงการ จึงไม่จำเป็นต้องติดตั้งฟอนต์เพิ่มในเครื่องผู้ใช้

\textbf{ตัวหนา} \textit{ตัวเอียง} และ English text can be used together.

\section{ทดสอบการตัดบรรทัดและกระจายข้อความแบบไทย}
\noindent
ภาษาไทยเขียนคำต่อเนื่องกันโดยทั่วไปและไม่ได้เว้นวรรคระหว่างคำเหมือนภาษาอังกฤษ
ดังนั้นระบบจัดหน้าเอกสารต้องรู้ตำแหน่งที่เหมาะสมสำหรับการตัดบรรทัด
ย่อหน้าทดสอบนี้จงใจเขียนให้มีความยาวหลายบรรทัดเพื่อยืนยันว่า XeLaTeX
สามารถตัดข้อความภาษาไทยตามขอบเขตของคำและกระจายช่องไฟให้ข้อความชิดขอบทั้งสองด้านได้
โดยไม่ทำให้ตัวอักษรซ้อนกันหรือหลุดออกนอกพื้นที่หน้ากระดาษ
เมื่อแก้ไขเนื้อหาในงานวิจัย ข้อความควรไหลต่อไปยังบรรทัดใหม่อย่างสม่ำเสมอ
และยังคงอ่านได้อย่างเป็นธรรมชาติแม้ในย่อหน้าที่มีคำภาษาอังกฤษ เช่น Docker และ LaTeX ปะปนอยู่ด้วย

\end{document}
